Trong khuôn khổ môn học Trí tuệ nhân tạo, đề tài đã xây dựng một quy trình hoàn chỉnh để phát hiện giao dịch thẻ tín dụng bất thường và nghi vấn gian lận, từ khâu nạp, làm sạch dữ liệu đến khâu thiết kế thử nghiệm, huấn luyện và đánh giá mô hình. Bộ dữ liệu gồm 30\,058 giao dịch trong giai đoạn 21/06--30/06/2020, trong đó 133 giao dịch gian lận, chiếm 0{,}44\%. Phân tích khám phá cho thấy gian lận tập trung ở vài loại dịch vụ (\texttt{shopping\_net}, \texttt{misc\_net}, \texttt{grocery\_pos} chiếm 82/133 vụ) và 104/133 vụ nằm trong vùng số tiền bất thường.

Đề tài xây dựng 14 đặc trưng, huấn luyện bốn họ thuật toán (Random Forest, XGBoost, LightGBM, CatBoost) với bốn chiến lược cân bằng mất cân bằng dữ liệu, kết hợp ba mô hình tốt nhất bằng stacking, chọn ngưỡng theo F2-Score trên tập Validation và đánh giá một lần duy nhất trên tập kiểm tra gồm 6\,012 giao dịch có 27 giao dịch gian lận.

Kết quả cho thấy kỹ thuật tạo đặc trưng đóng vai trò quyết định: với sáu đặc trưng gốc, PR-AUC chỉ đạt 0{,}5016, còn khi bổ sung nhóm đặc trưng kỹ thuật và mã hóa bằng \texttt{TargetEncoder} thì đạt 0{,}8197. So với \texttt{OneHotEncoder}, \texttt{TargetEncoder} vừa nâng cao hiệu quả vừa giảm số chiều từ 26 xuống 14 và rút ngắn thời gian tinh chỉnh còn 232 giây; ba thử nghiệm loại bỏ đặc trưng đều không làm mô hình tốt hơn. Xét tổng thể, LightGBM là thuật toán bền vững nhất: dẫn đầu PR-AUC ở cả năm thử nghiệm có nó và giữ FPR ở mức 0{,}0003--0{,}0017, còn các thuật toán khác biến động mạnh theo tập đặc trưng.

Vì vậy, đề tài chọn LightGBM với \texttt{scale\_pos\_weight} trên 14 đặc trưng mã hóa bằng \texttt{TargetEncoder} (thử nghiệm 3) làm mô hình đề xuất. Mô hình đạt PR-AUC 0{,}8197, F1-Score 0{,}8163 và F2-Score 0{,}7692 --- đều cao nhất trong sáu thử nghiệm --- với Precision 0{,}9091, Recall 0{,}7407; thời gian dự đoán toàn bộ 6\,012 giao dịch chỉ 0{,}036 giây.

Đề tài còn hạn chế do dữ liệu mô phỏng chỉ quan sát trong 10 ngày, số mẫu gian lận nhỏ, đặc trưng khoảng cách nhiễu tọa độ và xác suất dự đoán chưa hiệu chuẩn. Hướng phát triển là mở rộng dữ liệu thực tế, bổ sung đặc trưng hành vi thiết bị, hiệu chuẩn xác suất và chọn ngưỡng theo chi phí kinh doanh, rồi giám sát giao dịch. Tóm lại, đề tài đã minh họa quy trình phát hiện gian lận tài chính bằng trí tuệ nhân tạo và cung cấp cơ sở so sánh về hiệu quả các thuật toán. Với 80\% giao dịch hợp lệ không bị chặn nhầm và thời gian dự đoán dưới 0{,}04 giây, mô hình đề xuất đủ để thử nghiệm trong hệ thống chấm điểm rủi ro thời gian thực.
